%!TEX TS-program = xelatex
\documentclass[12pt, a4paper]{article}  

\usepackage{etex} % расширение классического tex в частности позволяет подгружать гораздо больше пакетов, чем мы и займёмся далее

%%%%%%%%%% Математика %%%%%%%%%%
\usepackage{amsmath,amsfonts,amssymb,amsthm,mathtools} 
%\mathtoolsset{showonlyrefs=true}  % Показывать номера только у тех формул, на которые есть \eqref{} в тексте.
%\usepackage{leqno} % Нумерация формул слева


%%%%%%%%%%%%%%%%%%%%%%%% Шрифты %%%%%%%%%%%%%%%%%%%%%%%%%%%%%%%%%
\usepackage{fontspec}         % пакет для подгрузки шрифтов
\setmainfont{Arial}   % задаёт основной шрифт документа

\defaultfontfeatures{Mapping=tex-text}

% why do we need \newfontfamily:
% http://tex.stackexchange.com/questions/91507/
\newfontfamily{\cyrillicfonttt}{Arial}
\newfontfamily{\cyrillicfont}{Arial}
\newfontfamily{\cyrillicfontsf}{Arial}

\usepackage{unicode-math}     % пакет для установки математического шрифта
\setmathfont{Asana Math}      % шрифт для математики
% \setmathfont[math-style=ISO]{Asana Math}
% Можно делать смену начертания с помощью разных стилей

% Конкретный символ из конкретного шрифта
% \setmathfont[range=\int]{Neo Euler}

\usepackage{polyglossia}      % Пакет, который позволяет подгружать русские буквы
\setdefaultlanguage{russian}  % Основной язык документа
\setotherlanguage{english}    % Второстепенный язык документа


%%%%%%%%%% Работа с картинками %%%%%%%%%
\usepackage{graphicx}                  % Для вставки рисунков
\usepackage{graphics}
\graphicspath{{images/}{pictures/}}    % можно указать папки с картинками
\usepackage{wrapfig}                   % Обтекание рисунков и таблиц текстом


%%%%%%%%%% Работа с таблицами %%%%%%%%%%
\usepackage{tabularx}            % новые типы колонок
\usepackage{tabulary}            % и ещё новые типы колонок
\usepackage{array}               % Дополнительная работа с таблицами
\usepackage{longtable}           % Длинные таблицы
\usepackage{multirow}            % Слияние строк в таблице
\usepackage{float}               % возможность позиционировать объекты в нужном месте
\usepackage{booktabs}            % таблицы как в книгах!
\renewcommand{\arraystretch}{1.3} % больше расстояние между строками

% Заповеди из документации к booktabs:
% 1. Будь проще! Глазам должно быть комфортно
% 2. Не используйте вертикальные линни
% 3. Не используйте двойные линии. Как правило, достаточно трёх горизонтальных линий
% 4. Единицы измерения - в шапку таблицы
% 5. Не сокращайте .1 вместо 0.1
% 6. Повторяющееся значение повторяйте, а не говорите "то же"
% 7. Есть сомнения? Выравнивай по левому краю!





%%%%%%%%%% Графика и рисование %%%%%%%%%%
\usepackage{tikz, pgfplots}  % язык для рисования графики из latex'a
\usepackage{amscd}                  %Пакеты для рисования
\usepackage[matrix,arrow,curve]{xy} %комунитативных диаграмм


%%%%%%%%%% Теоремы %%%%%%%%%%
\theoremstyle{plain}              % Это стиль по умолчанию.  Есть другие стили. 
\newtheorem{theorem}{Теорема}[section]
\newtheorem{result}{Следствие}[theorem]
% счётчик подчиняется теоремному, нумерация идёт по главам согласованно между собой

\theoremstyle{definition}         % убирает курсив и что-то еще наверное делает ;)
\newtheorem*{defin}{Определение}  % нумерация не идёт вообще

\newtheorem{fignia}{Какая-то фигня}




%%%%%%%%%% Свои команды %%%%%%%%%%
\usepackage{etoolbox}    % логические операторы для своих макросов


% Все свои команды лучше всего определять не по ходу текста, как это сделано в этом документе, а в преамбуле!




% Пакет, который ставит в каждом первом абзаце главы красную строку
% Просто, чтобы эта pdf-ка нормально смотрелась :)
\usepackage{indentfirst}
\setkeys{russian}{babelshorthands=true}


\title{Свои команды и макросы}
\date{\today}

\begin{document}

\maketitle


\section{Новые команды}
\subsection{Создание простых команд}

% Уже делали вот такие команды! 
% \DeclareMathOperator{\Var}{Var} можно только в прембуле!

\def \a{\alpha}
\def \R{\ensuremath{\mathbb{R}}} 

% \newcomand{имя команды}{Что делать} 

\newcommand{\RR}{\ensuremath{\mathbb{R}}}

Неважно как написать \R или \RR. По факту это будет одно и то же.

% ensuremath - при упоминании команды \R по ходу текста сам включает мтематический режим.

% Собственноручная расстановка пробелов
\newcommand{\iid}{i.\hspace{1.5pt}i.\hspace{1.5pt}d.}

$i.i.d.$ или $\iid$

\subsection{Команды с аргументами}
% \newcomand{имя команды}[количество аргументов]{Что делать}

\newcommand{\bb}[1]{\ensuremath{\mathbb{#1}} }

\bb{R} или \bb{Z} или \bb{Q}


$4/7$ или $^4/_7$

\newcommand{\fr}[2]{^#1/_#2}

$\fr{4}{7}$


\subsection{Переопределение команд}

$ \phi $ и $ \epsilon $

\renewcommand{\phi}{\varphi}
\renewcommand{\epsilon}{\varepsilon}

$ \phi $ и $ \epsilon $

\section{Счётчики}

% Можно объявить счётчик вот так:
% \newcounter{имя счётчика}[то чему подчиняется счётчик (не обязательно)]

\newcounter{jtem}[section]

% Счётчик можно установить на любую величину:
\setcounter{jtem}{13}

% Cчётчик можно вывести в любой из форм!
\arabic{jtem}

\alph{jtem}

\Alph{jtem}

\asbuk{jtem}  %(команда в babel или polyglossia) 

\Asbuk{jtem}

\roman{jtem}

\Roman{jtem}

% Можно выяснить какой номер у текущей главы или страницы
\arabic{section}

\arabic{page}

\Roman{page}

% вывести счётчики можно ещё и вот так
\thesection --- \thepage

% Можно поменять тип счётчика в нумерации!
% Сделаем номера страниц римскими
\renewcommand{\thepage}{\roman{page}}

% Сделаем номера глав буевенными
\renewcommand{\thesection}{\Asbuk{section}}
% При этом начнём нумерацию с буквы А
\setcounter{section}{0}

% Обратите внимание на то, что для глав счётчик сброшен, а для страниц нет!

% Дефолтные счётчики в LaTeX:  part, chapter, section, subsection, subsubsection, paragraph, subparagraph, page, figure, table, footnote, mpfootnote, equation
% Также определены счётчики для списков: enumi, enumii, enumiii, enumiv! 



\section{Списки}

\begin{enumerate}
	\renewcommand{\labelenumi}{\Roman{enumi}}
	\renewcommand{\labelenumii}{\asbuk{enumii}}
\item Первый пункт
\item Второй пункт
\item Третий пункт
\begin{enumerate}
\item Первый подпункт
\item Второй подпункт
\end{enumerate}
\end{enumerate}


\begin{enumerate}
\item Первый элемент
  \begin{enumerate}
  \item Первый вложенный элемент
  \item Второй вложенный элемент
  \end{enumerate}
\end{enumerate}

% Внутри списка другие не видят этого
% Тут или в преамбуле - видит только текущий список
\renewcommand{\labelenumi}{\arabic{enumi}}
\renewcommand{\labelenumii}{\arabic{enumi}.\arabic{enumii}}

\begin{enumerate}
\item Первый элемент
  \begin{enumerate}
  \item Первый вложенный элемент
  \item Второй вложенный элемент
  \end{enumerate}
\end{enumerate}



\section{Задача}

% объявляем новый счётчик!
\newcounter{i}
% задаём команду для оформления задач
\newcommand{\ex}[1]{%
\addtocounter{i}{1}    % увеличение счетчика на единицу
\textcolor{blue}{Задача \arabic{i}\\}
#1\\
}


\ex{В корзине лежат 12 яблок и 10 апельсинов. Ваня выбирает из неё яблоко или апельсин, после чего Надя берёт и яблоко, и апельсин. В каком случае Надя имеет большую свободу выбора: если Ваня взял яблоко или если он взял апельсин?}

\ex{Шестеро друзей пришли в театр! В ложе 6 мест. Среди друзей две девушки, которые не могут сидеть скраю. Сколько способов рассадки существует?}



\section{Теоремы}
\begin{defin}
\textbf{ Биноминальным коэффициентом} называется выражение $C_n^k = \frac{n!}{k!(n-k)!}$
\end{defin}

\begin{theorem}[Комбинаторное тождество]\label{th:1}
 \[C_n^0 + C_n^1 + C_n^2 + C_n^3 + \ldots + C_n^n = 2^n \]
\end{theorem}
\begin{proof}
Из Бинома Ньютона очевидно, что 

\[ 2^n = (1 + 1)^n = \sum_{k=0}^n C_n^k \cdot 1^k \cdot 1^{n-k} = \sum_{k=0}^n C_n^k. \] 
\end{proof}

\begin{result} \label{sl:11} 
Мощность множества всех подмножеств непустого множества равна $2^n$! 
\end{result}

Следствие \ref{sl:11} следует из теоремы \ref{th:1} со страницы \pageref{th:1}.


\begin{fignia}
	 Петя любит Катю, а Катя любит программировать и ботать по ночам! 
\end{fignia}


\section{Свои окружения} 
\subsection{Задача} 

% \newenvironment{имя}[число аргументов]{что должно быть до}{что должно быть после} 

% новый счётчик, ежели он нужен
\newcounter{iexer}

\newenvironment{exercise}
{ % то, что будет перед:
\addtocounter{iexer}{1}  % счётчик +1
 { 
   \noindent\ignorespaces   % игнор пробела
   \fontspec{Phorssa} % иной шрифт
   \textcolor{blue}{ \large  \textbf{Задача \arabic{iexer}}} % синий большой жирный текст 
 } 
 
\vspace{2mm} % расстояние между заголовком и текстом

%\noindent \ignorespaces \color{green} \fontsize{15}{1.2ex}\selectfont
% можно убрать лишние пробелы, поменять размер, цвет и тп для текста внутри окружения
} % конец того, что будет перед
{ % то, что будет после:
\begin{center}
\begin{tikzpicture}
\draw[blue,fill=blue](2,2) circle [radius = 0.3];
\end{tikzpicture}
\end{center}
} % конец того, что будет после

\begin{exercise}
В корзине лежат 12 яблок и 10 апельсинов. Ваня выбирает из неё яблоко или апельсин, после чего Надя берёт и яблоко, и апельсин. В каком случае Надя имеет большую свободу выбора: если Ваня взял яблоко или если он взял апельсин?
\end{exercise}


\subsection{Коробка} 

\newenvironment{korob}
    {\renewcommand{\arraystretch}{1}
    \begin{center}
    \begin{tabular}{|p{0.9\textwidth}|}
    \hline \vspace{1pt} 
    }{ 
    \vspace{5pt} \\ \hline 
    \end{tabular} 
    \end{center}
    }

 
\begin{korob}
Этот текст будет расположен внутри короба! Коробы это довольно интересные штуки.  
\end{korob}


% Можно придумывать для окружений аргументы!
\newcounter{ikorob}
\newenvironment{korob2}[1]
    {\renewcommand{\arraystretch}{1}
    \begin{center}
    \addtocounter{ikorob}{1} 
    Короб  \arabic{ikorob} --- #1 \\[1ex]
    \begin{tabular}{|p{0.9\textwidth}|}
    \hline \vspace{1pt} 
    }{ 
    \vspace{5pt} \\ \hline 
    \end{tabular} 
    \end{center}
    }
 
\begin{korob2}{Демонстрация короба}
Этот текст будет расположен внутри короба! Коробы это довольно интересные штуки.  
\end{korob2}



\section{etoolbox}
\subsection{Скрываем кусок текста} 

% Одно из применений - уничтожение какого-то куска текста!
\newbool{answers}
\booltrue{answers}
%\boolfalse{answers}

\renewcommand{\ex}[2][Нет ответа.]{%

\addtocounter{i}{1}    % увеличение счетчика на единицу
\textcolor{blue}{Задача \arabic{i}\\}
#2 

\ifbool{answers}{Ответ: #1}{ }
% {если} {правда} {ложь} 
}




\ex{В корзине лежат 12 яблок и 10 апельсинов. Ваня выбирает из неё яблоко или апельсин, после чего Надя берёт и яблоко, и апельсин. В каком случае Надя имеет большую свободу выбора: если Ваня взял яблоко или если он взял апельсин?}


\ex[много]{Шестеро друзей пришли в театр! В ложе 6 мест. Среди друзей две девушки, которые не могут сидеть скраю. Сколько способов рассадки существует?}



\subsection{Проверяем наличие аргумента} 

\NewDocumentCommand{\alice}{ m g }
 {%
  \IfNoValueTF{#2}%
   {%
    \it \noindent\ignorespaces #1 \\[5pt]
   \rightline{{Кто-то (Льюис Кэррол) }}
   }
   {
    \it \noindent\ignorespaces #1 \\[5pt]
    \rightline{{#2 (Льюис Кэррол) }}
   }%
 }

\alice{Видала я котов без улыбок, но улыбка без кота! Такого я в жизни ещё не встречала!}{Алиса} 

\alice{Видала я котов без улыбок, но улыбка без кота! Такого я в жизни ещё не встречала!} 



\end{document}

